\scanpage{57}
\mspage{50}

avec les localiseurs $W(I')$, $W(I)$, donc ils ne \deleted{peuvent} doivent pas être confondus avec les foncteurs $u_{\unsure{!|*}}^W,u_*^W$. \deleted{\unsure{précédent}} Donc j'ai été peut-être hâtif, dans ce que j'ai dit au sujet de $u_*$ et de $u_!$. Je serai obligé d'étudier la situation de très près, quand j'essayerai de construire les foncteurs $u_*^W$ et $u_!^W$. Pour le moment, je n'ai encore aucune assurance que les conditions envisagées \unsure{sur|pour} $(M,W)$ \inserted{suffisent à} \unsure{assurer|assurent} l'existence des $u_*^W,u_!^W$ -- pas $=$ pour $I,I'$ des ens. ordonnés finis !

% margin note written sideways

\margin{Comparer avec le lemme p. \unsure{63|13|15} et sa démonstration (plus exactement, le cor. du lemme).}

% slanted

\margin{\uline{Prop. 3} Soit $X_{\unsure{\bullet|.}}$ dans $M(\Delta^1)$, $X=(X_0\xleftarrow{\alpha}X_1)$. Alors $X_{\bullet}$ est dans \deleted{\unsure{$\mathrm{ob}.$}} $M(\Delta^1)_f$ \unsure{sss|ssi} $X_0,X_1$ dans $M_f$ et $\alpha\in F$.}

\margin{\deleted{\uline{\unsure{Proposition 3}}}\inserted{\unsure{Cor. à la|Cor. à 4|Cor. 1}} Soit $X\in M(I)_f$, alors pour \unsure{tte} $\alpha:s\to t$, dans $I$ qui est \uline{mono}, $X(\alpha):X(\unsure{t|f})\to X(s)$ est dans \uline{$F$}, (On applique cor. \uline{1} \inserted{(p. 47)} à $\widetilde{\alpha}:\Delta^1\to I$)}

\margin{[\uline{Cor.} ? Soit $\unsure{I'|II'}\overset{u}{\leftrightarrow}I$ \uline{imm. ouverte}. Alors \unsure{pour} $X\in\mathrm{ob}\,M(I)_f$, $\varprojlim_I X(i)\to\varprojlim_{I'}\unsure{\alpha|X}(i')$ est dans $F$.}

\textcircled{5}. \uline{Catégories de modèles quillénisables}.

Je reviens au cas général d'une catégorie de modèles $(M,W)$.

\llap{\raisebox{-1.75\baselineskip}[0pt][0pt]{\rule{0.4pt}{1.85\baselineskip}}\hspace{0.5em}}\uline{Définition 1} On dit que $(M,W)$ est \uline{quillénisable} s'il existe une structure de $\unsure{\mathrm{Qf}|\mathrm{QF}}$ \uline{exacte} $(C,\mathrm{TC},F,\mathrm{TF})$ pour $M$.

C'est là une condition autoduale, mais que je ne sais vérifier en pratique que moyennant l'exactitude soit à gauche, soit à droite -- donc par une condition plus forte qui détruit la symétrie. C'est pour la sauvegarder dans la mesure du possible que j'introduis la définition précédente, ainsi que la

\uline{Définition 2} Soit, en plus de $(M,W)$, $\deleted{\mathrm{Diag}}\subset\mathrm{Cat}$ une sous-catégorie pleine de $\mathrm{Cat}$. On dit que
% archive sheet number: 56
